\section{Incidence reconstruction and covariance}
\label{nuclear:5}
\subsection{The complete isometry and its domain}\label{nuc05:isometria}

Let \(H=\mathbb R^{12}\), with the Euclidean inner product; let \(P_0=I-\mathbf1\mathbf1^*/12\), \(\mu(k)=\mathbf1^*k/12\), and let \(\mathcal I:H\to\mathbb R^{132}\) be the incidence matrix of \cref{exc:entrelazador,exc:isometria}. Write \(E_{\rm inc}=\mathcal I(\mathbf1^\perp)\) and

\[
Y=\mathbb R\oplus E_{\rm inc},\qquad
\langle(\mu,y),(\nu,z)\rangle_Y=12\mu\nu+\langle y,z\rangle.
\]

The complete isometry constructed in Section~\ref{nuclear:2} is

\[
F:H\longrightarrow Y,\qquad Fk=(\mu(k),\mathcal I P_0k/6),
\]

\[
F^{-1}(\mu,y)=\mu\mathbf1+P_0\mathcal I^*y/6.
\]

The Gram identity \(\mathcal I^*\mathcal I=36I+30\mathbf1\mathbf1^*\) proves both inverse compositions and \(F^*F=I_H\), \(FF^*=I_Y\), with the indicated weighted inner product. Hence \(F\) is orthogonal between \(H\) and \(Y\). It can be complexified without changing the proofs.

The codomain is the twelve-dimensional space \(Y\), not the whole of \(\mathbb R\oplus\mathbb R^{132}\). Extending an operator from \(Y\) to its complement in that ambient space requires an additional choice; this is not part of the uniquely transported dynamics. In normalized signed coordinates, one uses \(FJ_H\), with \(J_H=H_{12}/2\). Integral inversion remains \(H_{12}/4\) on its sublattice, rather than an alternative metric normalizer.

\subsection{Covariant lifting theorem on the image}\label{nuc05:levantamiento}

\begin{theorem}
For a bounded operator \(D\) on \(H\), there exists a unique operator on \(Y\) satisfying \(D_YF=FD\). It is

\[
D_Y=FDF^{-1}.
\]

The correspondence preserves sums, products, adjoints, norm, positivity, and identity. In particular, it maps unitaries to unitaries and self-adjoint observables to self-adjoint observables. For a unitary \(C\) and an observable \(A\),

\[
C^*AC=A\quad\Longleftrightarrow\quad C_Y^*A_YC_Y=A_Y.
\]

If \(R\) is a unitary chart change, \(R_Y=FRF^{-1}\), then

\[
F(RDR^{-1})F^{-1}=R_YD_YR_Y^{-1}.
\]
\end{theorem}

\begin{proof}
The first formula follows by multiplying the intertwining relation by \(F^{-1}\). Uniqueness follows from surjectivity onto \(Y\). The product and adjoint identities follow from \(F^{-1}=F^*\), and equality of norms follows from the surjective isometry. The final equality is an associative composition of these maps. No commutation hypothesis between \(R\) and \(D\) enters.
\end{proof}

For \(g\in\operatorname{Aut}(\mathcal H)\), Proposition~\ref{exc:isometria} supplies

\[
F\rho_{12}(g)=\bigl(1\oplus\rho_{\rm hex}(g)\bigr)F.
\]

Thus, the chart change has a concrete combinatorial action on the hexads. For a general unitary \(D\), however, \(D_Y\) is an operator of the incidence representation, not necessarily a permutation of the design. Preserving the inner product of that representation does not make all its unitaries combinatorial automorphisms or elements of the Monster.

\subsubsection{Covariance under a noncommuting chart change}

Let \(C\) be the cyclic advance \(Ce_p=e_{p+1}\) on the twelve positions, namely the reverse direction of the convention \(S\) in \cref{sec:k-registro}. This is an operation of the window chart, not complete temporal holonomy. Let \(R\) be the already constructed star

\[
R=(1\ 3)(2\ 11)(5\ 7)(8\ 12).
\]

Then \(RCe_1=e_{11}\), whereas \(CRe_1=e_4\). They do not commute. Nevertheless, for \(C'=RCR^{-1}\), the identities \(T(C')R=RT(C)\) and \(\eta(C')R=R^{\oplus9}\eta(C)\) hold exactly, as do their versions on \(Y\) through \(F\). Chart covariance therefore works in an example where commutation in a fixed chart fails.

\subsection{Explicit recovery from the incidence record of complements}\label{nuc05:recuperacion}

Let \(C\) be unitary, \(T=(8I+C)/9\), \(\eta=(I-JJ^*)S_CJ:H\to H^9\), and let \(P\) be the projection onto \(\operatorname{Fix}(C)\), as in Section~\ref{nuclear:2}. Write \(F_9=F^{\oplus9}\). The incidence version of the finite record is

\[
\mathcal W_Nk=\left(FT^Nk,F_9\eta k,F_9\eta Tk,\ldots,F_9\eta T^{N-1}k\right).
\]

\begin{theorem}
This map is isometric and has the explicit left inverse

\[
\mathcal W_N^*(a,b_0,\ldots,b_{N-1})
=T^{*N}F^{-1}a+
\sum_{j=0}^{N-1}T^{*j}\eta^*F_9^{-1}b_j.
\]

In particular, applying this formula to \(\mathcal W_Nk\) recovers \(k\) exactly. The infinite record

\[
\mathcal W_\infty k=\left(FPk,(F_9\eta T^jk)_{j\ge0}\right)
\]

is also isometric, and its left inverse is

\[
k=PF^{-1}a+\sum_{j\ge0}T^{*j}\eta^*F_9^{-1}b_j
\quad\text{for }(a,(b_j))=\mathcal W_\infty k.
\]

The series converges in norm. For arbitrary square-summable data, the same formula defines the bounded adjoint, although not every datum in the codomain belongs to the image of \(\mathcal W_\infty\).
\end{theorem}

\begin{proof}
Applying \(F\) and \(F_9\) to each component of \(V_N\) or \(V_\infty\) preserves the inner product. Taking the adjoint of this composition gives the formulas. The telescopic identity in Section~\ref{nuclear:2} proves \(\mathcal W_N^*\mathcal W_N=I\), and its strong limit proves \(\mathcal W_\infty^*\mathcal W_\infty=I\). Convergence of the series follows by applying a bounded operator to truncations of a sequence in a Hilbert direct sum.
\end{proof}

This reconstruction does not preserve \(k\) as an intact first coordinate: the first component is the compressed terminal component or its fixed sector; the rest arises from successive complements. When \(P=0\), the incidence complements alone recover \(k\). The contribution recovered from the first \(N\) blocks is

\[
k_N=\sum_{j<N}T^{*j}\eta^*F_9^{-1}b_j
=\bigl(I-T^{*N}T^N\bigr)k.
\]

The remainder is \(T^{*N}T^Nk\), which tends to zero if \(P=0\). With a nonzero fixed sector, \(FPk\) must be retained; complement memory alone does not recover that sector. Reconstruction rates and dynamics on the block index are developed in Section~\ref{nuclear:3}.

If \(A\) is self-adjoint and \([A,C]=0\), the observable form is also preserved:

\[
\langle k,Ak\rangle
=\langle FPk,A_YFPk\rangle_Y
+\sum_{j\ge0}\langle F_9\eta T^jk,(I_9\otimes A_Y)F_9\eta T^jk\rangle.
\]

The scalar series converges absolutely because \(A\) is bounded and the sum of squared component norms is finite. The hypothesis expresses preservation of \(A\) by \(C\); it does not require \(C\) to commute with exceptional chart changes.

\subsection{Natural lifting of refinement with local frames}\label{nuc05:refinamiento}

The construction in \cref{iv:cont:historias,iv:cont:medida} supplies sets \(Z_n\) of enriched cylindrical states, truncations, and projective measures with full support. For a descendant \(z'\) of \(z\),

\[
w_n(z'|z)=\sqrt{\mu_{n+1}(z')/\mu_n(z)},\qquad
\sum_{\tau z'=z}w_n(z'|z)^2=1.
\]

The label \(z'\) preserves the ordered record, orientation, carry, boundary, memory, path, and terminal component. Refinement of a state does not identify distinct descendants.

Let \(q(z)\) be the integer memory record, \(C\) a unitary on \(H\), and \(R_z\) orthogonal changes of the fibre frame. For concrete incidence frames, one may use \(R_z=\rho_{12}(g_z)\), with \(g_z\in\operatorname{Aut}(\mathcal H)\), transporting their marks. Define

\[
G_z=R_zC^{q(z)},\qquad
U_{z',z}=G_{z'}G_z^*
=R_{z'}C^{q(z')-q(z)}R_z^*.
\]

This construction prolongs the explicit gauge equivalence of \cref{exc:doble-lectura}. It does not require \([R_z,C]=0\) and keeps the order of operations explicit.

On \(\mathscr H_n=\ell^2(Z_n)\otimes H\),

\[
\mathscr U_n(e_z\otimes v)
=\sum_{\tau z'=z}w_n(z'|z)e_{z'}\otimes U_{z',z}v.
\]

Define \(\mathscr F_n=I_{\ell^2(Z_n)}\otimes F\), with codomain \(\mathscr Y_n=\ell^2(Z_n)\otimes Y\).

\begin{theorem}
The refinement is isometric and satisfies

\[
\mathscr U^Y_n\mathscr F_n=\mathscr F_{n+1}\mathscr U_n,
\quad
\mathscr U^Y_n=\mathscr F_{n+1}\mathscr U_n\mathscr F_n^{-1}.
\]

The formula for each term is the preceding one with \(U^Y_{z',z}=FU_{z',z}F^{-1}\). Systems at several depths preserve composition and admit a surjective isometry between their inductive limits. All fibre coefficients are preserved exactly with their labels \(z\).
\end{theorem}

\begin{proof}
Descendants of distinct predecessors are orthogonal; each \(U_{z',z}\) is unitary; and the squared weights sum to one. This proves isometry. Along a trajectory, the intermediate factors \(G_z^*G_z\) cancel and the measure ratios telescope. Composition does not depend on the partition of the trajectory. The square with \(F\) is verified on \(e_z\otimes v\), term by term. The maps and their inverses are compatible isometries, so they extend as mutual inverses to the inductive completions.
\end{proof}

\subsubsection{Compatible observables and dynamics}

If \(A_0\) is a bounded fibre observable, let

\[
(\mathscr A_n\psi)_z=G_zA_0G_z^*\psi_z.
\]

Then \(\mathscr A_{n+1}\mathscr U_n=\mathscr U_n\mathscr A_n\), whence \(\mathscr U_n^*\mathscr A_{n+1}\mathscr U_n=\mathscr A_n\). This follows by substituting \(U_{z',z}=G_{z'}G_z^*\). It is conservation of an observable field transported by refinement; it does not assert that \(A_0\) commutes with \(C\) in a fixed chart.

The unitary dynamics at each level,

\[
(\mathscr C_n\psi)_z=G_zCG_z^*\psi_z
\]

satisfies \(\mathscr C_{n+1}\mathscr U_n=\mathscr U_n\mathscr C_n\). Their adjoints are also intertwined: since both \(\mathscr C\) are unitary, multiply the identity by their inverses. Thus, \(T\), its nonadic complements, and its fixed projections are intertwined. The finite and infinite records of Section~\ref{nuc05:recuperacion} are natural with respect to this refinement. Upon applying \(\mathscr F_n\), the same squares hold in the incidence fibres. For \(\mathscr A_n\) also to be conserved by these dynamics at a fixed level, the precise condition is \([A_0,C]=0\); covariance between frames still does not require \([R_z,C]=0\).

\subsubsection{Bilateral-memory specialization with varying frames}

On the reversible orbit of \cref{nuclear:3}, a coordinate readout may be represented by \(\psi=(\psi_m)_{m\in\mathbb Z}\in\ell^2(\mathbb Z;H)\). Its advance is \((C_0\psi)_m=\psi_{m-1}\). The map \(F\) acts pointwise. For frames \(R_m=\rho_{12}(g_m)\), the unitary change \((\mathscr R\psi)_m=R_m\psi_m\) produces

\[
(C'\psi)_m=R_mR_{m-1}^*\psi_{m-1}.
\]

The incidence operator has the same formula with \(1\oplus\rho_{\rm hex}(g_mg_{m-1}^{-1})\). Neither constant frames nor commutation with the shift is imposed. Since \(C'\) is conjugate to \(C_0\), it has no nonzero fixed vectors in this square-summable space. Section~\ref{nuc05:recuperacion} recovers the complete linear readout from its incidence complements. This concerns the represented orbit and coordinates, rather than an identification of all TPK histories with \(\mathbb Z\).

\subsection{Direct extension to the twelve \texorpdfstring{\(A_2\)}{A2} fibres}\label{nuc05:a2}

The exceptional construction preserves twelve \(A_2\) fibres and the order-three operator \(g_c\) developed in Section~\ref{nuc04:c3}. Let \(V_{A_2}=A_2\otimes_{\mathbb Z}\mathbb R\) be their two-dimensional Euclidean realization. The same positional isometry can be tensored over \(\mathbb R\) with this space:

\[
F_{A_2}=F\otimes I_{V_{A_2}}:
\mathbb R^{12}\otimes_{\mathbb R}V_{A_2}
\longrightarrow Y\otimes_{\mathbb R}V_{A_2}.
\]

Its first component is the vector mean, and its second consists of the centred incidence sums, now valued in \(V_{A_2}\). The incidence Gram operator, tensored with the identity, again proves the complete isometry and inverse. Hence \(g_c^Y=F_{A_2}g_cF_{A_2}^{-1}\) is orthogonal, of order three, and fixed-point-free. The nonadic composition of Section~\ref{nuclear:2} can be read out and recovered on this \(V_{A_2}\)-valued screen through Sections~\ref{nuc05:levantamiento} and~\ref{nuc05:recuperacion}. The twelve scalars of \(K\) are not confused with the twenty-four coordinates of the \(A_2\) fibres, nor is the action turned into a hexad permutation when the lift also acts within the fibres.

\subsection{Exact scope with respect to histories}

Let \(k:X\to H\) be the record reader on a history domain, and let \(\ell=F\circ k\). Then

\[
\ell(x)=\ell(y)\quad\Longleftrightarrow\quad k(x)=k(y).
\]

The complete screen adds zero loss relative to \(K\), but it does not change the fibres of this reader. A map \(d:X\to X\) descends to dynamics on the values of \(K\) if and only if

\[
k(x)=k(y)\Longrightarrow k(d x)=k(d y).
\]

The same condition is necessary and sufficient for descent to the incidence image. The proof consists of defining \(D(k(x))=k(dx)\); the condition is equivalent to independence of the representative. If this \(D\) is linear and unitary in the stated realization, Section~\ref{nuc05:levantamiento} applies. If it is nonlinear, \(F\) still conjugates it as a map but does not make it unitary.

The criterion distinguishes descent to \(K\) from the enriched dynamics. Section~\ref{subsec:k-terminal} establishes that the reader uses subpath, orientation, carry, and incidences; Section~\ref{exc:limite-inverso} distinguishes recovery of the visible word, support, and deep memory. On \(\ell^2(Z_n)\otimes H\), Section~\ref{nuc05:refinamiento} preserves enriched labels and transports the fibre representation. The equivalence acts on the previously generated history domain.

For an already constructed cut, naturality of the record in \cref{prop:k-extension} fixes \(K_N=K_{108}\circ\tau_{N,108}\): \(F\) preserves this naturality exactly. A sequence of new events may update \(K\); the additive or affine composition in \cref{sec:k-registro} preserves independence of partition, not constancy of the value when events are added.
